\chapter{Complex Vectors, Hermitian Matrices, and the Fourier Matrix}

Every length we have measured so far has been a real length. Every angle a real
angle. Our vectors lived in $\R^n$, our matrices had real entries, and the dot
product $\vec u\cdot\vec v=\sum u_iv_i$ answered every metric question we asked.
That setting is comfortable, and most of this book stays inside it.

But there is a quiet pressure building. Back in the eigenvalue chapters we kept
bumping into matrices -- innocent-looking real matrices -- whose eigenvalues
\emph{refused} to be real. A rotation in the plane has no real eigenvectors at
all; spin every arrow by ninety degrees and not a single direction survives.
Yet the characteristic polynomial still factors, just over a bigger number
system, and out come $\pm i$. The geometry was pushing us to enlarge the
number field, and this chapter is where we do it.

We are going to let our vectors have \emph{complex} entries. At first this looks
like a small bookkeeping change -- swap $\R$ for $\C$ -- but it forces one real
adjustment to the dot product, and that single adjustment ripples outward into
two new families of matrices: the \textbf{Hermitian} matrices (the complex cousins
of symmetric matrices) and the \textbf{unitary} matrices (the complex cousins of
orthogonal matrices). These are not exotic. They are the matrices of quantum
mechanics, of signal processing, of every algorithm that touches a frequency.

At the end we reach the most important unitary matrix in applied
mathematics: the \textbf{Fourier matrix}. It turns a signal
into its frequencies, and the engine underneath it -- the Fast Fourier Transform
-- is, by most accounts, the most-run numerical algorithm in the world.
Everything in this chapter is built to make that matrix feel inevitable.

Slow down at the very first idea, the conjugate in the inner product. Get that
one right and the rest of the chapter follows from it.

\chapterline

\section{Why real numbers are not enough}

Let us start with the rotation that broke things. The matrix
\[
  \mat R = \colvec{0 & -1 \\ 1 & \phantom{-}0}
\]
spins every vector in the plane by $90^\circ$ counterclockwise. Geometrically it
is obvious that \emph{no} nonzero real vector $\vec x$ satisfies $\mat R\vec x =
\lambda\vec x$: an eigenvector would have to point along the same line after
rotation, and a quarter turn sends every line to a different line. So over $\R$,
this matrix has no eigenvectors.

But the algebra does not give up. Its characteristic polynomial is
\[
  \det(\mat R - \lambda\mat I) = \det\colvec{-\lambda & -1 \\ 1 & -\lambda}
  = \lambda^2 + 1,
\]
and $\lambda^2+1=0$ has solutions $\lambda = \pm i$, where $i=\sqrt{-1}$. The
eigenvalues exist; they just live in $\C$. And once we allow complex entries in
the eigenvector too, genuine eigenvectors appear:
\[
  \mat R \colvec{1 \\ -i} = \colvec{0 & -1 \\ 1 & 0}\colvec{1\\-i}
  = \colvec{i \\ 1} = i\colvec{1 \\ -i},
\]
so $\colvec{1\\-i}$ is a genuine eigenvector with eigenvalue $i$. The geometry of
``no fixed direction'' and the algebra of ``complex eigenvalues'' are two faces of
the same fact. Enlarging to $\C$ does not destroy the geometry; it completes it.

\begin{intuitionbox}
A real matrix can ``rotate'' inside the plane spanned by two coordinates, and a
rotation has no real eigen-direction. Complex numbers are exactly the tool for
talking about rotation algebraically: multiplication by $e^{i\theta}$ \emph{is}
rotation by $\theta$. So when a real matrix wants to rotate, its eigenvalues come
out as complex numbers of the form $e^{i\theta}$, encoding the angle of the spin.
Complex eigenvalues are not a glitch -- they are the matrix telling you it turns.
\end{intuitionbox}

\subsection{A two-minute refresher on complex numbers}

We will need only a handful of facts; the appendix has the full story. A complex
number is $z = a + bi$ with $a,b\in\R$ and $i^2=-1$. We call $a=\operatorname{Re}z$
the real part and $b=\operatorname{Im}z$ the imaginary part. Two operations matter
enormously for us.

The \textbf{conjugate} flips the sign of the imaginary part:
\[
  \bar z = a - bi.
\]
The \textbf{modulus} (or absolute value) is the length of $z$ viewed as the point
$(a,b)$ in the plane:
\[
  \abs{z} = \sqrt{a^2+b^2}.
\]
The reason these two travel together is a single identity:
\[
  z\bar z = (a+bi)(a-bi) = a^2 - (bi)^2 = a^2 + b^2 = \abs{z}^2.
\]
Multiplying a number by its own conjugate \emph{kills the imaginary part} and
leaves a real, nonnegative number -- the squared length. That one fact drives
everything in this chapter.

\begin{figure}[H]
\centering
\begin{tikzpicture}[scale=1.2,>=Latex]
  \draw[->,gray!55] (-0.6,0) -- (3.2,0) node[right]{$\operatorname{Re}$};
  \draw[->,gray!55] (0,-1.8) -- (0,1.9) node[above]{$\operatorname{Im}$};
  \draw[->,very thick,vecblue] (0,0) -- (2.4,1.4) node[above right]{$z=a+bi$};
  \draw[->,very thick,vecred] (0,0) -- (2.4,-1.4) node[below right]{$\bar z=a-bi$};
  \draw[dashed,gray!60] (2.4,1.4) -- (2.4,-1.4);
  \node[vecblue] at (1.05,0.95) {$\abs{z}$};
  \draw[gray] (0.55,0) arc (0:30.3:0.55);
  \node at (0.85,0.18) {\small$\theta$};
  \draw[gray] (0.55,0) arc (0:-30.3:0.55);
  \fill[black] (0,0) circle (1pt);
\end{tikzpicture}
\caption{A complex number $z$ and its conjugate $\bar z$ are mirror images across
the real axis. Both have the same modulus $\abs z$; conjugation flips the angle
from $+\theta$ to $-\theta$. The product $z\bar z=\abs z^2$ is real.}
\label{ch32:conjugate}
\end{figure}

As Figure~\ref{ch32:conjugate} shows, $z$ and $\bar z$ are mirror images
across the real axis, sharing a modulus but carrying opposite angles.
The polar form makes the geometry vivid. Writing $z = r e^{i\theta}$ with
$r=\abs z$ and $\theta$ the angle, Euler's formula $e^{i\theta}=\cos\theta +
i\sin\theta$ says $z$ is the point at distance $r$ and angle $\theta$. Then
$\bar z = r e^{-i\theta}$ (same radius, opposite angle), and multiplication adds
angles: $(r_1e^{i\theta_1})(r_2e^{i\theta_2}) = r_1r_2\,e^{i(\theta_1+\theta_2)}$.
Multiplying by a number of modulus $1$ -- a number $e^{i\theta}$ on the unit
circle -- is a pure rotation. Keep that picture in mind for the Fourier section.

\section{Complex vectors and the right inner product}

Now stack complex numbers into a vector. A vector in $\C^n$ is just a list of $n$
complex numbers,
\[
  \vec z = \colvec{z_1 \\ z_2 \\ \vdots \\ z_n}, \qquad z_k\in\C.
\]
Addition and scalar multiplication work exactly as before, now with complex
scalars. The trouble -- and it is a productive trouble -- starts when we ask for
length.

\subsection{The naive dot product fails}

In $\R^n$ we defined $\norm{\vec u}^2 = \vec u\cdot\vec u = \sum u_i^2$, and the
whole edifice of geometry stood on one fact: $\vec u\cdot\vec u \ge 0$, zero only
for the zero vector. Try to copy it into $\C^n$ with the obvious formula
$\sum z_k^2$ and watch it collapse. Take the perfectly fine nonzero vector
$\vec z = \colvec{1 \\ i}$. Then
\[
  \sum z_k^2 = 1^2 + i^2 = 1 + (-1) = 0.
\]
A nonzero vector with ``length squared'' equal to zero. And it gets worse: the
vector $\colvec{1\\1+i}$ gives $1 + (1+i)^2 = 1 + 2i = $ a non-real number, so
``length squared'' is not even a real number we could take the square root of.
The naive copy is broken. We need a repair.

\begin{trapbox}
The number-one mistake when first meeting $\C^n$ is to write $\norm{\vec z}^2 =
\sum z_k^2$. Squaring a complex number does \emph{not} produce a length; $i^2=-1$,
not $+1$. Length must come from $\abs{z_k}^2 = z_k\bar z_k$, never from $z_k^2$.
The conjugate is not optional decoration -- it is the only thing standing between
you and ``nonzero vectors of length zero.''
\end{trapbox}

\subsection{The fix: conjugate one side}

We already saw the cure two sections ago. For a single number, the real,
nonnegative ``length squared'' is $z\bar z = \abs z^2$, not $z^2$. So for a vector
we should sum $z_k\bar z_k$:
\[
  \norm{\vec z}^2 \;=\; \sum_{k=1}^n z_k\bar z_k
  \;=\; \sum_{k=1}^n \abs{z_k}^2 \;\ge\; 0,
\]
and this is zero exactly when every $\abs{z_k}=0$, i.e. when $\vec z = \vec 0$.
Positivity is restored. Check our broken example: $\vec z=\colvec{1\\i}$ now has
$\norm{\vec z}^2 = \abs 1^2 + \abs i^2 = 1 + 1 = 2$, a sensible length
$\sqrt 2$.

Generalizing from $\vec z$ with itself to two different vectors gives the
\textbf{complex (or Hermitian) inner product}. We conjugate the entries of the
first vector and pair them with the second:
\[
  \boxed{\;\inner{\vec u}{\vec v} \;=\; \sum_{k=1}^n \bar u_k\, v_k
  \;=\; \bar u_1 v_1 + \bar u_2 v_2 + \cdots + \bar u_n v_n. \;}
\]
(Convention warning: physicists often conjugate the second slot instead. We pick
the first; the appendix notes the swap. Either way, exactly one side is
conjugated.) With this definition $\inner{\vec z}{\vec z} = \sum\bar z_kz_k =
\sum\abs{z_k}^2 = \norm{\vec z}^2$ is real and nonnegative -- the whole point.

\begin{definition}\label{ch32:inner}
For $\vec u,\vec v\in\C^n$, the \textbf{Hermitian inner product} is
$\inner{\vec u}{\vec v} = \sum_{k=1}^n \bar u_k v_k$, and the \textbf{norm} is
$\norm{\vec z} = \sqrt{\inner{\vec z}{\vec z}} = \sqrt{\sum_k\abs{z_k}^2}$.
Two vectors are \textbf{orthogonal} when $\inner{\vec u}{\vec v}=0$.
\end{definition}

\begin{example}\label{ch32:innerex}
Let $\vec u = \colvec{1+i \\ 2}$ and $\vec v = \colvec{3 \\ 1-i}$. Conjugate the
\emph{first} vector's entries: $\bar u_1 = 1-i$, $\bar u_2 = 2$. Then
\[
  \inner{\vec u}{\vec v} = (1-i)(3) + (2)(1-i) = (3-3i) + (2-2i) = 5 - 5i.
\]
And the squared length of $\vec u$ is
\[
  \norm{\vec u}^2 = \inner{\vec u}{\vec u}
  = (1-i)(1+i) + (2)(2) = (1+1) + 4 = 6,
\]
so $\norm{\vec u} = \sqrt 6$. The self inner product came out real and
positive even though $\vec u$ has complex entries -- the conjugate did its job.
\end{example}

\subsection{Conjugate symmetry: the price of the fix}

There is one property we must give up, and it is worth seeing exactly which one.
In $\R^n$ the dot product was symmetric: $\vec u\cdot\vec v = \vec v\cdot\vec u$.
In $\C^n$ symmetry becomes \textbf{conjugate symmetry}:
\[
  \inner{\vec v}{\vec u} \;=\; \overline{\inner{\vec u}{\vec v}}.
\]
Swapping the two vectors conjugates the answer. You can read it straight off the
definition: $\inner{\vec v}{\vec u} = \sum\bar v_k u_k$, and conjugating
$\inner{\vec u}{\vec v}=\sum\bar u_kv_k$ gives $\sum u_k\bar v_k =
\sum\bar v_k u_k$, the same thing. In Example~\ref{ch32:innerex} we found
$\inner{\vec u}{\vec v} = 5-5i$; you can check directly that
$\inner{\vec v}{\vec u} = (3)(1+i)+(1+i)(2) = 5+5i$, indeed the conjugate.

This is not a defect; it is the only way to keep $\inner{\vec z}{\vec z}$ real.
Conjugate symmetry \emph{forces} $\inner{\vec z}{\vec z} = \overline{\inner{\vec
z}{\vec z}}$, and a number equal to its own conjugate is exactly a real number.

\begin{formalbox}
The Hermitian inner product on $\C^n$ satisfies, for all $\vec u,\vec v,\vec w\in
\C^n$ and $c\in\C$:
\begin{enumerate}
  \item \textbf{Conjugate symmetry:} $\inner{\vec v}{\vec u} =
        \overline{\inner{\vec u}{\vec v}}$.
  \item \textbf{Linearity in the second slot:} $\inner{\vec u}{\vec v + \vec w} =
        \inner{\vec u}{\vec v} + \inner{\vec u}{\vec w}$ and
        $\inner{\vec u}{c\,\vec v} = c\,\inner{\vec u}{\vec v}$.
  \item \textbf{Conjugate-linearity in the first slot:}
        $\inner{c\,\vec u}{\vec v} = \bar c\,\inner{\vec u}{\vec v}$.
  \item \textbf{Positivity:} $\inner{\vec z}{\vec z} = \sum\abs{z_k}^2 \ge 0$, with
        equality iff $\vec z = \vec 0$.
\end{enumerate}
\end{formalbox}

Property 3 is the one that surprises people: a scalar pulled out of the
\emph{first} slot comes out conjugated, because that slot is the conjugated one.
If you scale the first vector by $i$, the inner product scales by $\bar i = -i$.
Everything else behaves just like the real dot product.

\begin{walkbox}
Take a walk and convince yourself of this without any formula: ``length must be a
real, nonnegative number, and it can only be zero for the zero vector.'' Now ask:
what is the \emph{simplest} change to $\sum u_kv_k$ that guarantees
$\inner{\vec z}{\vec z}\ge 0$ for every complex vector? You will rediscover the
conjugate. The mathematics had no choice; the conjugate is the unique repair.
\end{walkbox}

\section{The conjugate transpose}

We packaged the real dot product as a matrix product: $\vec u\cdot\vec v =
\vec u\T\vec v$, a row times a column. The complex inner product wants the same
tidy packaging, but the row has to carry that conjugate. So we bundle ``transpose
and conjugate'' into a single operation.

\begin{definition}\label{ch32:adjoint}
The \textbf{conjugate transpose} (or \textbf{Hermitian adjoint}) of a matrix
$\mat A\in\C^{m\times n}$ is $\mat A^{*}\in\C^{n\times m}$, obtained by transposing
$\mat A$ and conjugating every entry: $(\mat A^{*})_{jk} = \overline{(\mat A)_{kj}}$.
For a vector, $\vec u^{*}$ is the conjugated row.
\end{definition}

Many books write $\mat A^{H}$ (``$H$'' for Hermite) or $\mat A^{\dagger}$ (``dagger,''
the physics standard) for the same thing. We use the star $\mat A^{*}$. With it,
the inner product is just a product:
\[
  \inner{\vec u}{\vec v} \;=\; \vec u^{*}\vec v
  \;=\; \begin{bmatrix}\bar u_1 & \cdots & \bar u_n\end{bmatrix}
        \colvec{v_1\\\vdots\\v_n}.
\]
The star \emph{is} the inner product made into matrix algebra. Notice the clean
analogy: everywhere a real formula had a transpose $\mat A\T$, the complex version
has a conjugate transpose $\mat A^{*}$. The reversal rule survives unchanged,
$(\mat A\mat B)^{*} = \mat B^{*}\mat A^{*}$, and $(\mat A^{*})^{*} = \mat A$.

\begin{example}\label{ch32:starex}
\[
  \mat A = \colvec{1+i & 2 \\ 0 & 3-i}
  \quad\Longrightarrow\quad
  \mat A^{*} = \colvec{1-i & 0 \\ 2 & 3+i}.
\]
Transpose (swap across the diagonal), then conjugate every entry. The $2$ and $0$
are real, so they only move; the $1+i$ and $3-i$ both move \emph{and} flip the sign
of their imaginary part.
\end{example}

\begin{focusbox}
Translate ``real $\to$ complex'' by one rule: \textbf{every transpose becomes a
conjugate transpose.} Dot product $\vec u\T\vec v\to\vec u^{*}\vec v$. Symmetric
$\mat A\T=\mat A\to$ Hermitian $\mat A^{*}=\mat A$. Orthogonal
$\mat Q\T\mat Q=\mat I\to$ unitary $\mat U^{*}\mat U=\mat I$. The geometry is
identical; only the star changed.
\end{focusbox}

\section{Hermitian matrices: the complex symmetric}

A real symmetric matrix satisfied $\mat A\T=\mat A$ and gave us the crown jewel of
the real theory -- the Spectral Theorem: real eigenvalues, orthogonal
eigenvectors. We now ask which complex matrices deserve the same crown. The answer
is the matrices that equal their own conjugate transpose.

\begin{definition}\label{ch32:hermdef}
A square matrix $\mat A\in\C^{n\times n}$ is \textbf{Hermitian} if
$\mat A^{*} = \mat A$, that is, $a_{jk} = \overline{a_{kj}}$ for every entry.
\end{definition}

Two consequences are visible immediately from $a_{jk}=\overline{a_{kj}}$. On the
diagonal, $j=k$ forces $a_{kk}=\overline{a_{kk}}$, so \emph{the diagonal entries
are real}. Off the diagonal, entries come in conjugate-mirror pairs across the
diagonal. A real symmetric matrix is the special case where the off-diagonal pairs
are equal real numbers.

\begin{example}\label{ch32:hermex}
\[
  \mat A = \colvec{2 & 3-i \\ 3+i & 5}
\]
is Hermitian: the diagonal $2,5$ are real, and the off-diagonal entries $3-i$ and
$3+i$ are conjugates of each other. Transpose-and-conjugate and you get $\mat A$
back. By contrast $\colvec{2 & 3-i\\3-i & 5}$ is \emph{not} Hermitian -- the
off-diagonal pair must be conjugates, not equal.
\end{example}

\subsection{Real eigenvalues, for free}

The headline property is the same one symmetric matrices enjoyed, and the proof is
a two-line argument that shows exactly why we needed the conjugate inner
product.

\begin{theorem}\label{ch32:realeig}
Every eigenvalue of a Hermitian matrix is real.
\end{theorem}

\begin{proof}
Let $\mat A^{*}=\mat A$ and suppose $\mat A\vec x = \lambda\vec x$ with
$\vec x\ne\vec 0$. Form the number $\vec x^{*}\mat A\vec x$ and compute it two
ways. First, using $\mat A\vec x=\lambda\vec x$,
\[
  \vec x^{*}\mat A\vec x = \vec x^{*}(\lambda\vec x)
  = \lambda\,\vec x^{*}\vec x = \lambda\,\norm{\vec x}^2.
\]
Second, take the conjugate transpose of the scalar $\vec x^{*}\mat A\vec x$. A
$1\times 1$ matrix is its own transpose, so applying the star just conjugates the
number, while on the right it distributes by the reversal rule and uses
$\mat A^{*}=\mat A$:
\[
  \overline{\vec x^{*}\mat A\vec x} = (\vec x^{*}\mat A\vec x)^{*}
  = \vec x^{*}\mat A^{*}\vec x = \vec x^{*}\mat A\vec x.
\]
So $\vec x^{*}\mat A\vec x$ equals its own conjugate, hence is real. But it also
equals $\lambda\norm{\vec x}^2$, and $\norm{\vec x}^2>0$ is a positive real. A real
number divided by a positive real is real, so $\lambda$ is real.
\end{proof}

That argument is worth rereading. It is the complex twin of the symmetric proof,
and the conjugate inner product is doing all the work: $\norm{\vec x}^2$ is real
and positive \emph{because} we built the inner product to make it so. Swap in the
naive $\sum z_k^2$ and the proof falls apart -- another sign that the conjugate was
the right move.

\begin{aestheticbox}
There is something almost moral about this theorem. Hermitian matrices are the
matrices that ``measure'': in quantum mechanics every observable quantity --
energy, position, momentum, spin -- is a Hermitian operator, and
Theorem~\ref{ch32:realeig} is the reason a measurement can only return a
\emph{real} number. The mathematics guarantees, before any experiment, that the
possible outcomes (the eigenvalues) are real. The universe's insistence on real
measurements and the algebra's insistence on $\mat A^{*}=\mat A$ are the same
statement.
\end{aestheticbox}

\subsection{Orthogonal eigenvectors and the complex spectral theorem}

Just as in the real case, eigenvectors of a Hermitian matrix for \emph{different}
eigenvalues are orthogonal. The argument is again two lines. If $\mat A\vec x =
\lambda\vec x$ and $\mat A\vec y = \mu\vec y$ with $\lambda\ne\mu$ (both real, by
the theorem), then
\[
  \lambda\inner{\vec x}{\vec y} = \inner{\lambda\vec x}{\vec y}
  = \inner{\mat A\vec x}{\vec y} = \vec x^{*}\mat A^{*}\vec y
  = \vec x^{*}\mat A\vec y = \inner{\vec x}{\mat A\vec y}
  = \inner{\vec x}{\mu\vec y} = \mu\inner{\vec x}{\vec y}.
\]
(We used $\lambda$ real so it comes out of the conjugate slot unchanged, and
$\mat A^{*}=\mat A$ in the middle.) Thus $(\lambda-\mu)\inner{\vec x}{\vec y}=0$,
and since $\lambda\ne\mu$ we get $\inner{\vec x}{\vec y}=0$. Distinct eigenvalues
$\Rightarrow$ orthogonal eigenvectors, exactly as before.

Putting the pieces together gives the complex Spectral Theorem, the destination of
all this machinery.

\begin{theorem}[Spectral Theorem, complex form]\label{ch32:spectral}
Every Hermitian matrix $\mat A\in\C^{n\times n}$ can be written
\[
  \mat A = \mat U\,\mat\Lambda\,\mat U^{*},
\]
where $\mat\Lambda = \diag(\lambda_1,\dots,\lambda_n)$ has the (real) eigenvalues
on its diagonal and $\mat U$ is unitary, its columns an orthonormal basis of
eigenvectors. Real eigenvalues, orthonormal eigenvectors -- the same picture as
the real symmetric case, now over $\C$.
\end{theorem}

We will not grind out the full induction (it mirrors the real proof exactly), but
the moral is clean: \textbf{Hermitian is to $\C$ as symmetric is to $\R$.} Every
good thing symmetric matrices did, Hermitian matrices do, with stars in place of
transposes.

\begin{strangbox}
Strang teaches the real Spectral Theorem $\mat A=\mat Q\mat\Lambda\mat Q\T$ as the
summit of his eigenvalue chapters, then notes -- almost in passing -- that complex
entries simply replace $\T$ with the conjugate transpose. Reading this chapter
\emph{before} Strang, you will recognize his $\mat A=\mat Q\Lambda\mat Q\T$ as the
real shadow of $\mat A=\mat U\Lambda\mat U^{*}$. Reading it \emph{after}, you now
have the version that actually runs the Fourier transform and quantum mechanics --
the one his symmetric ``$\mathbf{S}=\mathbf{Q}\Lambda\mathbf{Q}\T$'' was secretly
pointing at.
\end{strangbox}

\section{Unitary matrices: the complex orthogonal}

The other star of the real theory was the orthogonal matrix $\mat Q$ with
$\mat Q\T\mat Q=\mat I$ -- the rotations and reflections, the rigid motions that
preserve every length and angle. Their complex analog is defined by the same
equation with a star.

\begin{definition}\label{ch32:unitarydef}
A square matrix $\mat U\in\C^{n\times n}$ is \textbf{unitary} if
$\mat U^{*}\mat U = \mat I$, equivalently $\mat U^{-1} = \mat U^{*}$. Its columns
form an orthonormal set in $\C^n$.
\end{definition}

Why ``orthonormal columns''? Read off the $(j,k)$ entry of $\mat U^{*}\mat U$: it
is (row $j$ of $\mat U^{*}$) times (column $k$ of $\mat U$) $=
\inner{\vec u_j}{\vec u_k}$, the inner product of two columns of $\mat U$.
Setting $\mat U^{*}\mat U=\mat I$ says $\inner{\vec u_j}{\vec u_k} = \delta_{jk}$:
each column has length $1$, and different columns are orthogonal. Identical to the
real story, with the Hermitian inner product.

\subsection{Unitary means length-preserving}

The defining geometric property -- the reason these matrices matter -- is that they
preserve the complex inner product, hence all lengths and angles.

\begin{theorem}\label{ch32:unitarynorm}
If $\mat U$ is unitary then $\inner{\mat U\vec x}{\mat U\vec y} =
\inner{\vec x}{\vec y}$ for all $\vec x,\vec y$, and in particular
$\norm{\mat U\vec x} = \norm{\vec x}$.
\end{theorem}

\begin{proof}
Slide the star across:
\[
  \inner{\mat U\vec x}{\mat U\vec y} = (\mat U\vec x)^{*}(\mat U\vec y)
  = \vec x^{*}\mat U^{*}\mat U\vec y = \vec x^{*}\mat I\vec y
  = \vec x^{*}\vec y = \inner{\vec x}{\vec y}.
\]
Setting $\vec y=\vec x$ gives $\norm{\mat U\vec x}^2 = \norm{\vec x}^2$.
\end{proof}

So a unitary matrix is a \textbf{rigid motion of $\C^n$}. It can rotate and reflect
and mix the coordinates as wildly as it likes, but it never stretches, never
shrinks, never distorts a single length. This is precisely why unitary matrices
are the ``safe'' transformations of numerical computing and of quantum mechanics:
applying one cannot amplify errors (it preserves their size) and cannot leak
probability (it preserves total length). And it explains the next fact in one
breath.

\begin{corollary}\label{ch32:unitaryeig}
Every eigenvalue $\lambda$ of a unitary matrix has $\abs\lambda = 1$: the
eigenvalues live on the unit circle.
\end{corollary}

\begin{proof}
If $\mat U\vec x=\lambda\vec x$ with $\vec x\ne\vec 0$, then
$\norm{\vec x} = \norm{\mat U\vec x} = \norm{\lambda\vec x} =
\abs\lambda\,\norm{\vec x}$, so $\abs\lambda = 1$.
\end{proof}

That is the perfect complement to the Hermitian story. Hermitian eigenvalues are
\emph{real} (they sit on the real line); unitary eigenvalues have \emph{modulus
one} (they sit on the unit circle). Each is the complex shadow of its real cousin:
symmetric eigenvalues are real, and orthogonal eigenvalues are $\pm 1$ (the
modulus-one numbers that happen to be real). Figure~\ref{ch32:eiglocus} puts the
two loci side by side.

\begin{figure}[H]
\centering
\begin{tikzpicture}[scale=1.25,>=Latex]
  % Hermitian: real line
  \begin{scope}[xshift=-3cm]
    \draw[->,gray!55] (-1.4,0) -- (1.4,0) node[right]{$\operatorname{Re}$};
    \draw[->,gray!55] (0,-1.3) -- (0,1.3) node[above]{$\operatorname{Im}$};
    \draw[very thick,vecblue] (-1.2,0) -- (1.2,0);
    \fill[vecblue] (-0.8,0) circle (1.6pt);
    \fill[vecblue] (0.3,0) circle (1.6pt);
    \fill[vecblue] (0.95,0) circle (1.6pt);
    \node[vecblue] at (0,-1.0) {\small Hermitian: real};
  \end{scope}
  % Unitary: unit circle
  \begin{scope}[xshift=2cm]
    \draw[->,gray!55] (-1.4,0) -- (1.4,0) node[right]{$\operatorname{Re}$};
    \draw[->,gray!55] (0,-1.3) -- (0,1.3) node[above]{$\operatorname{Im}$};
    \draw[very thick,vecred] (0,0) circle (1);
    \fill[vecred] (0.81,0.59) circle (1.6pt);
    \fill[vecred] (-0.31,0.95) circle (1.6pt);
    \fill[vecred] (-0.81,-0.59) circle (1.6pt);
    \node[vecred] at (0,-1.35) {\small unitary: $\abs\lambda=1$};
  \end{scope}
\end{tikzpicture}
\caption{Where the eigenvalues live. Hermitian matrices have real eigenvalues (on
the real axis); unitary matrices have eigenvalues of modulus one (on the unit
circle). These are the complex versions of ``symmetric $\Rightarrow$ real
eigenvalues'' and ``orthogonal $\Rightarrow$ eigenvalues $\pm 1$.''}
\label{ch32:eiglocus}
\end{figure}

\begin{example}\label{ch32:unitaryex}
The matrix
\[
  \mat U = \frac{1}{\sqrt 2}\colvec{1 & 1 \\ i & -i}
\]
is unitary. Its columns are $\vec u_1=\tfrac1{\sqrt2}\colvec{1\\i}$ and
$\vec u_2=\tfrac1{\sqrt2}\colvec{1\\-i}$. Check:
$\inner{\vec u_1}{\vec u_1} = \tfrac12(\bar 1\cdot 1 + \bar i\cdot i)
= \tfrac12(1 + (-i)(i)) = \tfrac12(1+1)=1$, length one; and
$\inner{\vec u_1}{\vec u_2} = \tfrac12(\bar 1\cdot 1 + \bar i\cdot(-i))
= \tfrac12(1 + (-i)(-i)) = \tfrac12(1 + i^2) = \tfrac12(1-1) = 0$, orthogonal.
Two orthonormal columns, so $\mat U^{*}\mat U=\mat I$. (Watch the conjugate on the
$i$ in the first slot -- without it you would wrongly get $\inner{\vec u_1}{\vec
u_1}=\tfrac12(1+i^2)=0$, the same trap as before.)
\end{example}

\begin{connectionbox}
Unitary matrices are the gateway to \emph{group theory}. The set of all $n\times n$
unitary matrices is closed under multiplication and inverse -- a product of rigid
motions is a rigid motion -- so it forms a group, written $U(n)$, the
\emph{unitary group}. Its members of determinant $1$ form $SU(n)$, and $SU(2)$,
$SU(3)$ are the symmetry groups at the heart of the Standard Model of particle
physics. The single equation $\mat U^{*}\mat U=\mat I$ in this chapter is the front
door to the symmetry language of modern physics. The real analog
$\mat Q\T\mat Q=\mat I$ gives the rotation groups $O(n)$ and $SO(n)$ -- the
symmetries of spheres and rigid bodies.
\end{connectionbox}

\section{The discrete Fourier matrix}

Now we build the most useful unitary matrix in applied mathematics. The idea is
an old one: a complicated signal -- a sound wave, a stock price, an image
row -- can be rebuilt as a \emph{sum of pure oscillations} of different
frequencies. The Fourier matrix is the change of basis that swaps the ``samples''
description of a signal for its ``frequencies'' description. What makes it work
is that this change of basis is \emph{unitary} -- a rigid rotation of $\C^n$ -- so
it is fully reversible and stable.

\subsection{Roots of unity: $n$ points evenly on the circle}

Everything is built from one number. Fix a size $n$ and define
\[
  \omega = e^{2\pi i/n} = \cos\frac{2\pi}{n} + i\sin\frac{2\pi}{n}.
\]
This $\omega$ is the point on the unit circle at angle $2\pi/n$ -- one $n$-th of the
way around. Its powers $1,\omega,\omega^2,\dots,\omega^{n-1}$ are $n$ points spaced
evenly around the unit circle, the \textbf{$n$-th roots of unity}: each is a
solution of $z^n=1$. They close up perfectly, $\omega^n = e^{2\pi i}=1$, so the
powers cycle with period $n$ ($\omega^{n+1}=\omega$, and so on).
Figure~\ref{ch32:roots} draws the eight evenly spaced roots for $n=8$.

\begin{figure}[H]
\centering
\begin{tikzpicture}[scale=1.6,>=Latex]
  \draw[gray!55] (0,0) circle (1);
  \draw[->,gray!55] (-1.3,0) -- (1.3,0) node[right]{$\operatorname{Re}$};
  \draw[->,gray!55] (0,-1.3) -- (0,1.3) node[above]{$\operatorname{Im}$};
  \foreach \k in {0,1,2,3,4,5,6,7} {
    \pgfmathsetmacro\ang{45*\k}
    \fill[vecblue] (\ang:1) circle (1.4pt);
    \draw[vecblue!40] (0,0) -- (\ang:1);
  }
  \node[vecblue,below right] at (0:1) {$\omega^0{=}1$};
  \node[vecblue,above right] at (45:1) {$\omega$};
  \node[vecblue,above] at (90:1) {$\omega^2$};
  \node[vecblue,left] at (180:1) {$\omega^4$};
  \draw[->,vecred,thick] (8:0.42) arc (8:37:0.42);
  \node[vecred] at (22:0.62) {\small$\tfrac{2\pi}{n}$};
\end{tikzpicture}
\caption{The eighth roots of unity ($n=8$): the powers of $\omega=e^{2\pi i/8}$ step
around the unit circle in equal jumps of $2\pi/n$. After $n$ steps you return to
$1$. These rotating numbers are the entire content of the Fourier matrix.}
\label{ch32:roots}
\end{figure}

The single fact we will lean on is a geometric one: \textbf{the $n$ roots of unity
sum to zero} (for $n\ge 2$). Picture $n$ equally spaced arrows from the center of
the circle -- by symmetry they pull equally in all directions and cancel. (Two
arrows: $1+(-1)=0$. Three: an equilateral triangle of arrows summing to the
center.) Algebraically, $1+\omega+\dots+\omega^{n-1}$ is a geometric series equal
to $\frac{\omega^n - 1}{\omega-1} = \frac{1-1}{\omega-1}=0$, since $\omega^n=1$ and
$\omega\ne 1$. Remember this cancellation; it is the gear that makes the Fourier
matrix unitary.

\subsection{Building the matrix}

The \textbf{$n\times n$ Fourier matrix} $\mat F = \mat F_n$ has, in row $j$ and
column $k$ (indexing from $0$), the entry
\[
  (\mat F)_{jk} = \omega^{\,jk}, \qquad j,k = 0,1,\dots,n-1.
\]
Row $j$ is the list of powers $1,\omega^j,\omega^{2j},\dots,\omega^{(n-1)j}$ -- a
single pure oscillation sampled at $n$ points, oscillating $j$ times around the
circle. Multiplying a signal $\vec x$ by $\mat F$ asks, in each row $j$, ``how much
of frequency $j$ is in $\vec x$?'' For $n=4$, with $\omega = e^{2\pi i/4} = i$
(since a quarter turn lands on $i$), the powers cycle $i^0,i^1,i^2,i^3 =
1,i,-1,-i$, and
\[
  \mat F_4 = \colvec{
    1 & 1 & 1 & 1 \\
    1 & i & -1 & -i \\
    1 & -1 & 1 & -1 \\
    1 & -i & -1 & i }.
\]
Each entry is $i^{jk}$ reduced using $i^4=1$. The top row is all $1$'s (frequency
zero -- it just sums the signal, the average/DC term). The second row $1,i,-1,-i$
is one full rotation; the bottom row is the same rotation run the other way.

\begin{labbox}
Compute $\mat F_4\,\vec x$ by hand for $\vec x = \colvec{1\\0\\0\\0}$ (a single
spike at time $0$) and then for $\vec x=\colvec{1\\1\\1\\1}$ (a constant signal).
You should find $\mat F_4\colvec{1\\0\\0\\0}=\colvec{1\\1\\1\\1}$ -- a spike
contains \emph{every} frequency equally -- and $\mat F_4\colvec{1\\1\\1\\1}=
\colvec{4\\0\\0\\0}$ -- a constant signal is \emph{pure} frequency zero. A spike in
time spreads across all frequencies; a flat line in time is a spike in frequency.
That trade-off is the seed of the uncertainty principle.
\end{labbox}

To see why each row deserves to be called ``a frequency,'' picture its real part
$\cos(2\pi jm/n)$ as $m$ runs over the $n$ sample points. Row $0$ is flat (no
oscillation -- the DC term), row $1$ completes one cycle across the samples, row
$2$ completes two, and so on. Figure~\ref{ch32:fourierrows} traces these sampled
cosines for $n=8$: the Fourier matrix is just ``all the pure oscillations that fit
on $n$ points,'' stacked as rows.

\begin{figure}[H]
\centering
\begin{tikzpicture}[>=Latex,xscale=0.62,yscale=0.85]
  \foreach \j/\yoff/\lbl in {0/0/{row $0$: flat (DC)}, 1/{-2.2}/{row $1$: one cycle}, 2/{-4.4}/{row $2$: two cycles}, 3/{-6.6}/{row $3$: three cycles}} {
    % axis line
    \draw[gray!45] (-0.3,\yoff) -- (8.3,\yoff);
    \node[anchor=east,font=\small] at (-0.6,\yoff) {\lbl};
    % continuous cosine
    \draw[vecblue,thick,domain=0:8,smooth,samples=120,variable=\t]
      plot ({\t},{\yoff + 0.75*cos(deg(2*pi*\j*\t/8))});
    % sampled stems at m=0..7
    \foreach \m in {0,1,2,3,4,5,6,7} {
      \pgfmathsetmacro\yv{0.75*cos(deg(2*pi*\j*\m/8))}
      \fill[vecred] (\m,\yoff+\yv) circle (1.6pt);
      \draw[vecred!55] (\m,\yoff) -- (\m,\yoff+\yv);
    }
  }
\end{tikzpicture}
\caption{The real parts of the first four rows of $\mat F_8$, read as oscillations
sampled (red dots) at the eight points $m=0,\dots,7$. Row $j$ completes exactly $j$
full cycles across the samples, so higher rows mean higher frequency. The Fourier
matrix bundles every such pure oscillation that fits on $n$ points, and asks how
much of each is present in a signal.}
\label{ch32:fourierrows}
\end{figure}

\subsection{Why $\mat F$ is (almost) unitary}

Here is the payoff. The columns of $\mat F$ are \emph{orthogonal} under the
Hermitian inner product. Take column $j$ and column $k$; their inner product
conjugates the first (and $\overline{\omega^{mj}} = \omega^{-mj}$, since
$\bar\omega = \omega^{-1}$ for a unit-circle number):
\[
  \inner{\text{col }j}{\text{col }k}
  = \sum_{m=0}^{n-1}\overline{\omega^{mj}}\,\omega^{mk}
  = \sum_{m=0}^{n-1}\omega^{m(k-j)}
  = \sum_{m=0}^{n-1}\bigl(\omega^{\,k-j}\bigr)^{m}.
\]
This is a geometric series in $\omega^{k-j}$. If $j=k$ every term is
$\omega^0 = 1$, so the sum is $n$. If $j\ne k$ then $\omega^{k-j}$ is itself a
nontrivial root of unity, and by our cancellation fact the sum is $0$. In one line:
\[
  \inner{\text{col }j}{\text{col }k} = \begin{cases} n & j=k,\\[2pt] 0 & j\ne k.
  \end{cases}
\]
Different columns are orthogonal; each column has squared length $n$. So
$\mat F^{*}\mat F = n\mat I$. The columns are orthogonal but have length $\sqrt n$,
not $1$ -- so $\mat F$ itself is not quite unitary. Divide by $\sqrt n$ to fix it.

\begin{theorem}\label{ch32:fourierunitary}
The \textbf{normalized Fourier matrix} $\;\mat U = \dfrac{1}{\sqrt n}\,\mat F\;$ is
unitary: $\mat U^{*}\mat U = \mat I$, equivalently $\mat F^{*}\mat F = n\mat I$ and
$\mat F^{-1} = \tfrac1n\mat F^{*}$.
\end{theorem}

This is the structural heart of the whole subject. The Fourier transform is a
\emph{rotation of $\C^n$}. It does not distort the signal -- it merely views the
same vector in a new orthonormal basis, the basis of pure frequencies. And because
the inverse of a unitary matrix is just its conjugate transpose, the inverse
Fourier transform costs no more than the forward one:
\[
  \mat F^{-1} = \frac1n\,\mat F^{*}, \qquad\text{so}\qquad
  \vec x = \frac1n\,\mat F^{*}(\mat F\vec x).
\]
You can always get your signal back. Going to the frequency world and coming home
are the same operation up to a conjugate and a factor of $n$.

\begin{example}\label{ch32:invex}
For $n=4$ we have $\mat F_4^{*}=\overline{\mat F_4\T}$. Since $\mat F_4$ is
symmetric ($\omega^{jk}=\omega^{kj}$), transposing changes nothing, and
conjugating sends $i\to-i$:
\[
  \mat F_4^{*} = \colvec{
    1 & 1 & 1 & 1 \\
    1 & -i & -1 & i \\
    1 & -1 & 1 & -1 \\
    1 & i & -1 & -i }, \qquad
  \tfrac14\mat F_4^{*}\mat F_4 = \mat I.
\]
You can verify a single entry: row $0$ of $\mat F_4^{*}$ dotted with column $0$ of
$\mat F_4$ is $1+1+1+1=4$, and divided by $4$ gives the $1$ on the diagonal of
$\mat I$. Row $0$ with column $1$ is $1+i-1-i=0$ -- the roots-of-unity
cancellation, doing its job.
\end{example}

\begin{trapbox}
Conventions for the Fourier matrix differ across textbooks and software, and they
will bite you. Some define $(\mat F)_{jk}=\omega^{-jk}$ (minus sign in the
exponent), some fold the $\tfrac1{\sqrt n}$ into the forward transform, some into
the inverse, some split it as $\tfrac1{\sqrt n}$ on each. All are ``the Fourier
matrix.'' The invariant facts that survive every convention: the columns are
orthogonal roots-of-unity vectors, and the matrix is unitary after the right
normalization. Always check which convention your library (NumPy, MATLAB) uses
before comparing numbers.
\end{trapbox}

\section{From the matrix to the FFT}

We have a beautiful object; now make it fast. Naively, multiplying $\mat F_n$ by a
vector $\vec x$ is an $n\times n$ matrix times an $n$-vector: about $n^2$ complex
multiplications. For a single second of CD-quality audio, $n\approx 44{,}000$, so
$n^2\approx 2$ billion multiplications -- per second of sound, for one transform.
That was, for a long time, simply too slow to be practical.

The \textbf{Fast Fourier Transform} (FFT) computes the very same $\mat F_n\vec x$
in about $n\log_2 n$ operations instead of $n^2$. For $n=44{,}000$ that is the
difference between $2$ billion and roughly $700{,}000$ -- a speedup of thousands.
The idea, in one sentence, is \emph{divide and conquer using the structure of the
roots of unity}.

\subsection{The one identity that makes it work}

Split a length-$n$ signal (with $n$ even) into its even-indexed and odd-indexed
entries: $\vec x_{\text{even}} = (x_0,x_2,x_4,\dots)$ and $\vec x_{\text{odd}} =
(x_1,x_3,x_5,\dots)$, each of length $n/2$. The miracle is that the length-$n$
transform of $\vec x$ can be assembled from two length-$n/2$ transforms -- one of
the evens, one of the odds. Writing $\omega_n=e^{2\pi i/n}$ and using
$\omega_n^2=\omega_{n/2}$ (going twice as fast around the circle is the same as one
step around a half-size circle), one can show
\[
  (\mat F_n\vec x)_j \;=\; E_j \;+\; \omega_n^{\,j}\,O_j,
  \qquad
  (\mat F_n\vec x)_{j+n/2} \;=\; E_j \;-\; \omega_n^{\,j}\,O_j,
\]
for $j=0,\dots,n/2-1$, where $E_j=(\mat F_{n/2}\vec x_{\text{even}})_j$ and
$O_j=(\mat F_{n/2}\vec x_{\text{odd}})_j$ are the half-size transforms (read
periodically in $j$). The two outputs $j$ and $j+n/2$ share the very same
$E_j,O_j$ and differ only by a sign on the twiddle term $\omega_n^{\,j}O_j$ -- the
sign flip comes from $\omega_n^{\,j+n/2}=-\omega_n^{\,j}$, since stepping a
further half-circle negates a unit-circle number. That paired $+/-$ combine is the
famous FFT \emph{butterfly}, and it is why one combine step produces two outputs at
once. Two half-size problems plus a cheap combine step. Recurse: each half splits
into
quarters, and so on, $\log_2 n$ levels deep, each level costing $O(n)$ work to
combine. Total: $O(n\log n)$. That is the whole algorithm in spirit -- the roots of
unity were \emph{designed} to factor this way. Figure~\ref{ch32:fft} draws the
recursion tree for $n=8$.

\begin{figure}[H]
\centering
\begin{tikzpicture}[scale=1,>=Latex,
   fnode/.style={draw,rounded corners,fill=blue!5,minimum width=2.2cm,minimum height=0.7cm,font=\small}]
  \node[fnode] (top) at (0,2.4) {$\mat F_8$ on $\vec x$ \ ($n=8$)};
  \node[fnode] (l1) at (-2.6,1.2) {$\mat F_4$ evens};
  \node[fnode] (r1) at (2.6,1.2) {$\mat F_4$ odds};
  \node[fnode] (ll) at (-3.9,0) {$\mat F_2$};
  \node[fnode] (lr) at (-1.3,0) {$\mat F_2$};
  \node[fnode] (rl) at (1.3,0) {$\mat F_2$};
  \node[fnode] (rr) at (3.9,0) {$\mat F_2$};
  \draw[->] (top) -- (l1);  \draw[->] (top) -- (r1);
  \draw[->] (l1) -- (ll);   \draw[->] (l1) -- (lr);
  \draw[->] (r1) -- (rl);   \draw[->] (r1) -- (rr);
  \node[font=\small,gray] at (5.9,2.4) {1 split};
  \node[font=\small,gray] at (5.9,1.2) {2 splits};
  \node[font=\small,gray] at (5.9,0) {$\log_2 8=3$};
\end{tikzpicture}
\caption{Divide and conquer. One length-$8$ transform splits into two length-$4$
transforms (evens and odds), each into two length-$2$ transforms. With $\log_2 n$
levels and $O(n)$ combine work per level, the cost is $O(n\log n)$ instead of
$O(n^2)$. The structure of the roots of unity is what allows the split.}
\label{ch32:fft}
\end{figure}

\begin{applicationbox}
The FFT is, by most counts, the most-executed numerical algorithm on Earth, and it
is hiding everywhere:
\begin{itemize}
  \item \textbf{Audio and music.} MP3, AAC, and Shazam all transform sound into
        frequencies; compression throws away the frequencies your ear cannot hear,
        and recognition matches a song by its frequency fingerprint.
  \item \textbf{Images.} JPEG splits a photo into $8\times 8$ blocks and applies a
        cousin of the Fourier transform (the discrete cosine transform), then
        discards small high-frequency coefficients -- that is where the compression
        comes from.
  \item \textbf{Wireless.} Wi-Fi, 4G/5G, and DSL use OFDM, which sends data on many
        frequency channels at once; the receiver runs an FFT to pull them apart.
  \item \textbf{Science and ML.} MRI scanners literally measure the Fourier
        transform of your body and FFT it back to an image; convolutions in signal
        processing (and fast polynomial multiplication) become cheap pointwise
        products in the frequency domain, because the Fourier basis diagonalizes
        convolution.
\end{itemize}
Every one of these rides on two facts from this chapter: the transform is unitary
(so it is invertible and stable), and it is fast (the FFT).
\end{applicationbox}

\begin{connectionbox}
The discrete Fourier transform is the finite shadow of the \emph{Fourier series}
from analysis, where a periodic function $f(t)$ is written as an infinite sum
$\sum_k c_k e^{ikt}$ of pure waves. The functions $e^{ikt}$ are orthonormal under
the inner product $\inner{f}{g}=\frac1{2\pi}\int_0^{2\pi}\overline{f(t)}g(t)\,dt$ --
the \emph{same} conjugate-inner-product idea from earlier, now with an integral in
place of a sum. Sampling such a function at $n$ equally spaced points turns the
integral into the finite sum that is our matrix $\mat F_n$, and ``orthonormal
waves'' becomes ``unitary matrix.'' Hermitian inner products, orthogonal bases of
oscillations, completeness -- this is the doorway to \emph{Fourier analysis} and
\emph{functional analysis}, where vectors become functions and matrices become
operators.
\end{connectionbox}

\section{A unified picture}

Step back and look at the symmetry between the real world and the complex world we
just built. The translation is mechanical and total: replace every transpose by a
conjugate transpose, and every theorem survives.

\begin{center}
\begin{tabular}{@{}lll@{}}
\toprule
\textbf{Concept} & \textbf{Real ($\R^n$)} & \textbf{Complex ($\C^n$)}\\
\midrule
Inner product & $\vec u\T\vec v=\sum u_iv_i$ & $\vec u^{*}\vec v=\sum\bar u_iv_i$\\
``Transpose'' & $\mat A\T$ & $\mat A^{*}$ (conjugate transpose)\\
Self-adjoint & symmetric $\mat A\T=\mat A$ & Hermitian $\mat A^{*}=\mat A$\\
\quad eigenvalues & real & real\\
Length-preserving & orthogonal $\mat Q\T\mat Q=\mat I$ & unitary $\mat U^{*}\mat U=\mat I$\\
\quad eigenvalues & $\pm 1$ & on unit circle, $\abs\lambda=1$\\
Spectral theorem & $\mat A=\mat Q\Lambda\mat Q\T$ & $\mat A=\mat U\Lambda\mat U^{*}$\\
\bottomrule
\end{tabular}
\end{center}

\begin{focusbox}
If you remember one thing from this chapter, remember the dictionary above. Complex
linear algebra is not a new subject -- it is real linear algebra with the conjugate
inserted at exactly one spot (one side of every inner product), and with the
transpose upgraded to the conjugate transpose to match. Everything else, including
the Spectral Theorem and the geometry of rigid motions, carries over word for word.
\end{focusbox}

\begin{recoverybox}
\textbf{Cost and intuition recap.} The conjugate inner product
$\inner{\vec u}{\vec v}=\sum\bar u_kv_k$ costs the same $n$ multiplications as the
real dot product -- one extra conjugation per entry, no asymptotic change -- and it
is the one repair that keeps $\norm{\vec z}^2$ real and nonnegative. A Hermitian
matrix gives \emph{real} eigenvalues; a unitary matrix preserves length and gives
eigenvalues of modulus one. The dense Fourier matrix-vector product costs $O(n^2)$;
the FFT exploits the roots-of-unity structure to bring it down to $O(n\log n)$,
which is the whole reason the frequency domain is practical at scale. One mantra
carries the chapter: \emph{star replaces transpose, and the conjugate goes on
exactly one side.}
\end{recoverybox}

\begin{aestheticbox}
The Fourier matrix is where three strands of this entire book braid together. It is
a \emph{change of basis}: from samples to frequencies. It is a \emph{unitary
matrix}: a rigid rotation of $\C^n$ that loses no information. And it
\emph{diagonalizes} every circular convolution at once, so the hardest operation in
signal processing becomes pointwise multiplication. One matrix, built from $n$
points spinning evenly around a circle, sits at the intersection of geometry,
algebra, and computation -- and it runs, billions of times a second, inside the
device you are reading this on.
\end{aestheticbox}

\section{Exercises}

\begin{exercise}\label{ch32:ex-innerprod}
Let $\vec u=\colvec{2i\\1}$ and $\vec v=\colvec{1+i\\3i}$ in $\C^2$. Compute
$\inner{\vec u}{\vec v}$ and $\inner{\vec v}{\vec u}$, and verify that
$\inner{\vec v}{\vec u}=\overline{\inner{\vec u}{\vec v}}$. Then compute
$\norm{\vec u}$ and $\norm{\vec v}$.
\end{exercise}

\begin{exercise}\label{ch32:ex-naivefail}
Find a nonzero vector $\vec z\in\C^2$ with $\sum z_k^2 = 0$ but
$\norm{\vec z}^2=\sum\abs{z_k}^2 \ne 0$. Explain in one sentence why this shows the
naive ``dot product'' $\sum z_k^2$ cannot define a length.
\end{exercise}

\begin{exercise}\label{ch32:ex-star}
For $\mat A=\colvec{2 & i\\ -i & 4}$ and $\mat B=\colvec{1 & 1+i\\ 1-i & 0}$,
compute $\mat A^{*}$ and $\mat B^{*}$. Which of $\mat A,\mat B$ are Hermitian?
For the Hermitian one(s), confirm by inspection that the diagonal is real and the
off-diagonal entries are conjugates.
\end{exercise}

\begin{exercise}\label{ch32:ex-hermeig}
The matrix $\mat A=\colvec{3 & 2i\\ -2i & 3}$ is Hermitian. Find its eigenvalues
(they should come out real, by Theorem~\ref{ch32:realeig}) and a pair of orthogonal
eigenvectors. Verify the orthogonality with the Hermitian inner product.
\end{exercise}

\begin{exercise}\label{ch32:ex-unitarycheck}
Show that $\mat Q=\colvec{\cos\theta & -\sin\theta\\ \sin\theta & \cos\theta}$ is
unitary for every real $\theta$ (it is even orthogonal). Then show that
$\mat V=\tfrac1{\sqrt2}\colvec{1 & i\\ i & 1}$ is unitary, and find a vector whose
length it preserves to double-check Theorem~\ref{ch32:unitarynorm}.
\end{exercise}

\begin{exercise}\label{ch32:ex-rootsum}
Using $\omega=e^{2\pi i/6}$, list the six sixth-roots of unity in the form
$e^{2\pi ik/6}$ and plot them on the unit circle. Verify (by the geometric-series
argument, or numerically) that they sum to zero. What is $\omega^3$?
\end{exercise}

\begin{exercise}\label{ch32:ex-f4}
Using $\mat F_4$ from the text, compute $\mat F_4\vec x$ for
$\vec x=\colvec{1\\-1\\1\\-1}$. Which single frequency does the answer pick out, and
why does that match the fact that $\vec x$ alternates sign every step (the
\emph{highest} frequency representable on $4$ points)?
\end{exercise}

\begin{exercise}\label{ch32:ex-f4unitary}
Verify directly that $\tfrac14\mat F_4^{*}\mat F_4=\mat I$ by computing the inner
product of column $1$ with column $1$ (should be $4$) and column $1$ with column $3$
(should be $0$) of $\mat F_4$. Identify the roots-of-unity cancellation in the
second computation.
\end{exercise}

\begin{exercise}\label{ch32:ex-unitarysim}
\textbf{(Proof.)} Prove that unitary similarity preserves eigenvalues: if $\mat U$
is unitary and $\mat A=\mat U^{*}\mat A_0\mat U$, then $\mat A$ and $\mat A_0$ have
the same characteristic polynomial, hence the same eigenvalues. Then prove the
special case used in the chapter: if $\mat A_0$ is Hermitian, so is $\mat A$.
\end{exercise}

\begin{exercise}\label{ch32:ex-normalproof}
\textbf{(Challenge.)} A matrix is called \textbf{normal} if $\mat A^{*}\mat A=
\mat A\mat A^{*}$. Verify that every Hermitian matrix and every unitary matrix is
normal. Then explain, in a sentence or two, why ``normal'' is exactly the class of
matrices to which the complex Spectral Theorem applies (eigenvectors form an
orthonormal basis) -- making Hermitian and unitary two special cases of one bigger
truth.
\end{exercise}

\begin{exercise}\label{ch32:ex-fft-count}
\textbf{(FFT counting.)} Let $T(n)$ be the number of combine operations the FFT
uses on a signal of length $n=2^m$. The divide-and-conquer step gives roughly
$T(n)=2\,T(n/2)+n$ with $T(1)=0$. Unroll this recurrence to show $T(n)=n\log_2 n$,
and compare $T(1024)$ with the naive $1024^2$.
\end{exercise}
